% TOP_PROBLEM: {"id": 1356, "title": "Open Coloring Axiom and Continuum Hypothesis", "category_id": 10, "category": {"id": 10, "name": "set_theory", "display_name": "Set Theory", "description": "Foundations of mathematics, infinite sets, and cardinality.", "slug": "set-theory", "order_index": 10, "created_at": "2026-07-31T15:26:25.671Z"}, "status": "open"}
% FIRST_POSTED: 2026-10-01
% ATTEMPT_STATUS: unresolved
% Imported from: unsolvedmath_additional_500.tex; UnsolvedMath ID 1356
% !TeX program = lualatex
% Compile twice: lualatex 1356.tex
% Self-contained: no external bibliography, images, or \input files.
% Numeric section labels and visible numbers are original UnsolvedMath IDs.
\documentclass[11pt,a4paper]{article}
\usepackage[margin=24mm,headheight=14pt]{geometry}
\usepackage{fontspec}
\setmainfont{Latin Modern Roman}
\setsansfont{Latin Modern Sans}
\setmonofont{Latin Modern Mono}
\usepackage{amsmath,amssymb,mathtools}
\usepackage{unicode-math}
\setmathfont{Latin Modern Math}
\usepackage{microtype}
\usepackage{enumitem,xurl,needspace}
\usepackage[dvipsnames]{xcolor}
\usepackage{hyperref}
\hypersetup{unicode=true,colorlinks=true,linkcolor=MidnightBlue,urlcolor=MidnightBlue,
 pdftitle={UnsolvedMath Problem 1356},pdfauthor={Source-annotated compilation},bookmarksnumbered=true}
\usepackage{bookmark,tocloft,titlesec,fancyhdr}
\setlength{\cftsecnumwidth}{4.2em}
\cftsetpnumwidth{2.5em}
\cftsetrmarg{4em}
\titleformat{\section}{\Large\bfseries\raggedright}{\thesection}{0.7em}{}
\pagestyle{fancy}\fancyhf{}
\fancyhead[L]{\small\sffamily Additional UnsolvedMath statements}
\fancyhead[R]{\small Original catalogue IDs}
\fancyfoot[C]{\thepage}
\setlength{\parindent}{0pt}
\setlength{\parskip}{5pt plus 1pt}
\setlength{\emergencystretch}{3em}
\setcounter{tocdepth}{1}\setcounter{secnumdepth}{1}
\setlist[itemize]{leftmargin=3em,itemsep=4pt,topsep=4pt,parsep=0pt}
\allowdisplaybreaks
\newcommand{\N}{\mathbb{N}}
\newcommand{\Z}{\mathbb{Z}}
\newcommand{\Q}{\mathbb{Q}}
\newcommand{\R}{\mathbb{R}}
\newcommand{\C}{\mathbb{C}}
\newcommand{\F}{\mathbb{F}}
\newcommand{\A}{\mathbb{A}}
\newcommand{\PP}{\mathbb{P}}
\newcommand{\E}{\mathbb{E}}
\newcommand{\eps}{\varepsilon}
\newcommand{\dd}{\,\mathrm d}
\newcommand{\sref}[2]{\hyperlink{source:#1:#2}{[S#2]}}
\newif\ifshowcatalogue
\showcataloguetrue
\newcommand{\cataloguescope}[1]{\ifshowcatalogue
 \paragraph{Statement recorded in the supplied source.}
 #1\fi}
\begin{document}
% UnsolvedMath ID 1356; source catalogue code ST-010

\setcounter{section}{1355}

\section{OCA with a Continuum Larger Than Aleph-Two}\label{1356}

{\small\sffamily UnsolvedMath ID 1356 · Set Theory}

\subsection{Definitions and mathematical statement}

\paragraph{Shared notation.}Unless a section says otherwise, $\N=\{1,2,\ldots\}$,
$\Z$ is the ring of integers, $\Q,\R,\C$ are the rational, real, and complex
fields, $[N]=\{1,\ldots,\lfloor N\rfloor\}$, and $\log$ is the natural
logarithm. For real functions of a parameter tending to infinity,
$f\ll g$ and $f=O(g)$ mean $|f|\leq Cg$ eventually for some fixed $C$;
$f\gg g$ reverses this comparison; $f=o(g)$ means $f/g\to0$;
$f\sim g$ means $f/g\to1$; and $f\asymp g$ means both order bounds.
A subscript on $O$ or $\ll$ specifies allowed constant dependence.
The expression $n^{o(1)}$ in an upper bound means growth smaller than
$n^\epsilon$ for each fixed $\epsilon>0$. Local definitions govern any
exception, finite-field convention, or use of the same symbol for another
object. An asymptotic claim is not automatically an effective algorithm.



OCA asserts that for a separable metric space $X$ and an open set $K_0\subseteq[X]^2$, either an uncountable $Y\subseteq X$ has $[Y]^2\subseteq K_0$, or $X=\bigcup_{n<\omega}X_n$ with $[X_n]^2\cap K_0=\varnothing$ for every $n$. Openness is in the induced unordered-pair topology. Ask whether ZFC together with OCA and $2^{\aleph_0}>\aleph_2$ is consistent. \sref{1356}{1} \sref{1356}{2}

\paragraph{Statement recorded in the supplied source.}
Is the open coloring axiom (OCA) consistent with $2^{\aleph_0} > \aleph_2$? \sref{1356}{1}

\subsection{Short English statement}

Can the open colouring axiom hold in a set-theoretic universe where the continuum is strictly larger than aleph-two?

\subsection{Research attempt: proving the Polish-space case and identifying the missing completeness}
\textbf{Status and source audit.} No consistency model with a large continuum is constructed. The inspected arXiv record for [S2] is by Thomas Gilton and Itay Neeman, and its theorem concerns the Abraham--Rubin--Shelah axiom. The target here is Todorčević's open-graph dichotomy. The inspected primary paper [S3] explicitly separates these as OCA[ARS] and OCA[T]. Their similar names cannot be used to transfer the large-continuum result.

\subsubsection{1. Removing all open pieces with countable colourings}
We prove the target dichotomy in ZFC when $X$ is Polish, meaning separable and completely metrizable. Fix a compatible complete metric and an open graph $K\subseteq[X]^2$. A set is independent if none of its distinct pairs lies in $K$, and countably colourable if it is a countable union of independent sets.

Let $U$ be the union of all open countably colourable subsets of $X$. Because $X$ has a countable base, $U$ is already the union of countably many such basic open sets. Countable unions of countably colourable sets are countably colourable, so $U$ is countably colourable. If $U=X$, this is the second alternative of OCA.

Otherwise let $Y=X\setminus U$, a nonempty closed Polish subspace. Every nonempty relatively open subset $V$ of $Y$ fails to be countably colourable. Indeed write $V=O\cap Y$ with $O$ open in $X$. If $V$ were countably colourable, then $O=(O\cap U)\cup V$ would be so too, giving $O\subseteq U$ by definition, contradicting $V\ne\varnothing$.

\subsubsection{2. Building a perfect clique}
Every nonempty relatively open subset of $Y$ therefore contains a $K$-edge; otherwise it would itself be independent. Recursively choose nonempty relatively open sets $V_s\subseteq Y$ indexed by finite binary strings. Arrange that the closures of $V_{s0}$ and $V_{s1}$ lie inside $V_s$, are disjoint, have diameter at most $2^{-|s|-1}$, and have every cross pair in $K$.

Here is why each step is possible. Select an edge $\{x,y\}$ inside $V_s$. Openness of the graph off the diagonal provides product neighborhoods of $x,y$ whose cross pairs are all edges. Small enough metric balls have disjoint closures inside those neighborhoods and inside $V_s$. The property of $Y$ ensures the resulting open pieces remain available for the next step.

For each binary sequence $b\in2^{\mathbb N}$, the nested nonempty closed sets $\overline{V_{b\upharpoonright n}}$ have diameters tending to zero. Completeness gives a unique point $x_b$ in their intersection. At the first digit where $b$ and $b'$ differ, the cross-pair condition gives $\{x_b,x_{b'}\}\in K$, and disjoint closures ensure the points are distinct. The map $b\mapsto x_b$ is continuous by the diameter condition. Its compact image is a perfect set, and in particular an uncountable clique. This proves the first alternative.

\subsubsection{3. Why arbitrary separable metric spaces remain outside the proof}
In a noncomplete subspace the nested closed sets can converge to a point of its completion that is not in the original space. For example, neighborhoods of a fixed irrational point intersected with $\mathbb Q$ can be nonempty, nested and of shrinking diameter while their intersection in $\mathbb Q$ is empty. The Cantor construction would then produce points outside the domain of the colouring, and no clique in $X$ follows.

Passing to the completion does not fix this automatically. A graph open on $X$ has an open extension to the ambient pair space, but the perfect clique obtained there may consist largely or entirely of new points. OCA quantifies over arbitrary separable subspaces, not only closed or completely metrizable ones.

\subsubsection{4. Verification and exact residual target}
The proof audit checks the countable-base reduction for $U$, closedness and completeness of $Y$, and closure containment in the open cross-edge neighborhoods. Each is essential. The argument works for Polish spaces regardless of the value of the continuum and thus supplies no upper bound of $\aleph_2$ on it. Extending the dichotomy to every arbitrary separable subspace while preserving a continuum larger than $\aleph_2$ is the unresolved consistency problem. Neither the ZFC special case nor the separately named OCA[ARS] theorem establishes that extension.

\subsection{Sources}

{\small

\begin{itemize}

\item[\hypertarget{source:1356:1}{S1}] ULAM.AI, \emph{UnsolvedMath}, supplied \texttt{problems.json}, record 1356 (ST-010), statement and background. The embedded literature-review date is 2026-08-17; that is the archive’s review date, not a fresh certification. Dataset landing page: \url{https://huggingface.co/datasets/ulamai/UnsolvedMath}.

\item[\hypertarget{source:1356:2}{S2}] T. Gilton and I. Neeman, Abraham-Rubin-Shelah Open Colorings and a Large Continuum, arXiv:1904.10516. \url{https://arxiv.org/abs/1904.10516}. Bibliographic information imported from the supplied record.

\item[\hypertarget{source:1356:3}{S3}] OCA with large continuum open-problem discussion. \url{https://pi.math.cornell.edu/~justin/Ftp/OCA_c.pdf}. Bibliographic information imported from the supplied record.

\end{itemize}

}

\label{end:1356}
\end{document}
